% 03_body.tex — 산출물 3/5 이해관계자 등록부와 현저성 분석 (빈 템플릿 / 완성 예시 공용 본문, A4 가로)
% 드라이버: 03_이해관계자등록부_빈템플릿.tex (\wbblanktrue) / 03_이해관계자등록부_완성예시.tex (\wbblankfalse)
% 내용 출처: PM트랙/Day1_워크북.md §3.3(항목 가이드) §3.4(빈 템플릿) §3.5(온담 P1 완성 예시본)
% 등록부 16열은 인쇄 폭 때문에 ID로 연결되는 두 패널(A: 식별·관심사 / B: 현저성·관여)로 나눠 배치했다.
\documentclass[10pt,a4paper]{article}
\usepackage[a4paper,landscape,margin=14mm,top=12mm,bottom=18mm]{geometry}
\def\DocNum{3}
\def\DocTitle{이해관계자 등록부와 현저성 분석}
\def\DocSub{Stakeholder Register \& Salience Analysis}
\def\DocVariant{\B{빈 템플릿}{온담 P1 완성 예시}}
\usepackage{wbstyle}
\begin{document}

\docheader
 {\Bg{[정식 명칭과 코드]}{차세대 주문·재고 통합 플랫폼 구축 (ONE ONDAM P1)}}
 {\Bg{[v1.x / 판정 시점: 게이트·날짜]}{v1.0 / 판정 시점 G0 2026-01-05}}
 {\Bg{[PM 직책·실명]}{P1 PM\rd{7mm}}}
 {\Bg{[갱신 책임자 — RACI 26행 A]}{P1 PM (RACI \#26 A: 게이트마다 갱신)\rd{7mm}}}

% ------------------------------------------------------------------ A
\secbar[0mm]{A. 등록부 — 식별 정보와 관심사 (항목 1~3)}
\guide{개인은 실명·직책, 조직은 단위명과 접점자. 내부/외부와 계약 관계 유무(수행사·3PL·벤더 vs 가맹점주·노조·규제기관)를 구분. 역할은 헌장 항목 11·RACI와 일치. 기대·우려는 \textbf{그 사람의 말로}, 분리해서 적고 KPI·공식 입장을 붙인다. “프로젝트 성공”이라고만 적는 것은 분석이 아니다.}
{\scriptsize
\begin{longtable}{|L{9mm}|L{30mm}|L{19mm}|L{32mm}|L{19mm}|L{46mm}|L{46mm}|L{46mm}|}\hline
\hd{ID} & \hd{성명 · 직책 · 소속} & \hd{내부·외부 (계약 유무)} & \hd{프로젝트 역할} & \hd{유형} & \hd{기대} & \hd{우려} & \hd{KPI · 공식 입장} \\ \hline \endhead
\ifwbblank
\rd{12mm}S-01 & \g{[실명 · 직책 · 소속]} & \g{[내부 / 외부(계약·비계약)]} & \g{[Exec / SU / SS / CCB / 수행사 / 인수자 / 감리 / 없음]} & \g{[의사결정자 / 요구제공자 / 사용자 / 공급자 / 규제·감독 / 영향받는 자]} & \g{[그 사람의 말로]} & \g{[잃을까 두려워하는 것]} & \g{[KPI, 공식 발언]} \\ \hline
\rd{12mm}S-02 & & & & & & & \\ \hline
\rd{12mm}S-03 & & & & & & & \\ \hline
\rd{12mm}S-04 & & & & & & & \\ \hline
\rd{12mm}S-05 & & & & & & & \\ \hline
\rd{12mm}S-06 & & & & & & & \\ \hline
\rd{12mm}S-07 & & & & & & & \\ \hline
\rd{12mm}S-08 & & & & & & & \\ \hline
\rd{12mm}S-09 & & & & & & & \\ \hline
\rd{12mm}S-10 & & & & & & & \\ \hline
\rd{12mm}S-11 & & & & & & & \\ \hline
\rd{12mm}S-12 & & & & & & & \\ \hline
\rd{12mm}S-13 & & & & & & & \\ \hline
\rd{12mm}S-14 & & & & & & & \\ \hline
\rd{12mm}S-15 & & & & & & & \\ \hline
\rd{12mm}S-16 & & & & & & & \\ \hline
\else
S-01 & 정미란 · CFO & 내부 & Executive, CCB 위원장 & 의사결정자 & 편익 근거가 있는 집행, 부채비율 160\% 이하 & 1차연도 176억 집행, 근거 없는 편익 & “편익 숫자는 근거가 있어야 승인” \\ \hline
S-02 & 김선호 · 대표이사 & 내부 & 프로그램 이사회 의장 & 의사결정자 & IPO 예비심사 전 가동, “가능하면 더 앞당기라” & 일정 지연 & “일정이 먼저다” \\ \hline
S-03 & 오정근 · 영업본부장 & 내부 & Senior User (P1·P3), 편익 책임자, CCB & 의사결정자·요구 제공자 & 매장이 교육 없이 쓰는 시스템, 성수기 무중단 & 2019년 반복, 가맹 반발, 재계약률 하락 & 재계약률 91\% KPI. “가맹점주가 반발하면 2019년의 반복” \\ \hline
S-04 & 박세연 · IT본부장 & 내부 & Senior Supplier, 아키텍처 최종 승인 & 의사결정자·공급자 & 자체 운영 가능, 벤더 비종속 & 상용SW·클라우드 비중, OD-POS·문서 부재 13본이 “실패”로 서술되는 것 & “우리 손으로 운영할 수 있어야” \\ \hline
S-05 & 한동석 · 물류본부장 & 내부 & P2 Senior User, 3PL 계약 소관, 재고 정확도 책임 & 의사결정자·요구 제공자 & 물류 자동화 우선, WMS 연동은 대성과 협의 & 처리량 4,200 근거, 노조 & “4,200은 현장에서 나온 숫자가 아니다” \\ \hline
S-06 & 나영선 · 품질안전본부장 & 내부 & CCB 위원, 품질 데이터 요구 제공 & 요구 제공자 & 검사기록 디지털화, 품질 데이터 접근 & P1 범위에서 품질이 빠지는 것 & “재고 정확도보다 검사기록이 급하다” \\ \hline
S-07 & 최영주 · 준법지원팀장 CPO & 내부 & CCB 위원, 규제 요건 해석 & 규제·감독 & 시행일 준수, 처리항목 142개 반영 & 검토 의견이 결정으로 오해되는 것 & “검토 의견을 제공하지, 아키텍처를 결정하지 않는다” \\ \hline
S-08 & 박준영 · 서비스운영팀장 & 내부 & P5 인수자, SAC 판정권 & 인수자 & 운영 문서·SLA 완비 후 인수 & 오픈 첫 달 사고가 콜센터로 & “문서와 SLA 없이는 인수하지 않는다” \\ \hline
S-09 & 이재현 · 한빛디지털 PM & 외부(계약) & 수행사 PM, 변경협의체 대표 & 공급자 & FP 기준선 준수, 손익률 6.5\% & 발주사 미결정 대기, 무상 범위 증가 & “대기와 추가는 변경대가 대상” \\ \hline
S-10 & 서정필 · 가맹점주협의회장 & 외부(비계약 — 가맹 계약의 집단 대표) & 가맹 포털 사용자 대표, 부속서 협상 상대 & 사용자·영향받는 자 & 발주 편의, 리드타임 단축 & 판매 데이터 용도, 단말 비용, 앱 주문의 가맹 몫 & “발주는 환영, 데이터·단말 비용은 별개” \\ \hline
S-11 & 강윤지 · 직영 지역매니저 & 내부 & 현장 사용자, 슈퍼유저 후보 & 사용자·영향받는 자 & “이번엔 없어지지 않는” 시스템 & 시스템 전환 때마다 두 배 업무 & 2019 스마트오더 중단 경험 \\ \hline
S-12 & 윤태호 · 내부감사팀장 & 내부 & 관리 체계 서면 질의자 & 규제·감독 & 편익 근거·컨틴전시 승인 근거 문서 & 2020·2023 지적 반복 & “근거를 문서로 제출하십시오” \\ \hline
S-13 & 김미르 · 노조 물류지부장 & 내부(비계약 대표) & 이천센터 246명 대표 & 영향받는 자 & 처리량 목표의 인력 영향에 대한 답 & 인력 감축 & “처리량 두 배가 인원 절반이면 누가 승인했는가” \\ \hline
S-14 & 조성태 · 생산본부장 겸 온담푸드시스템 대표 & 내부/별도 법인 & 공장 데이터 연동의 계약 주체 & 의사결정자(별도 법인) & 공장 출하 재고 연동의 실익 & 별도 법인 계약·보안 정책 부담 & 서면 등장 \\ \hline
S-15 & 강현도 · 구매팀장 & 내부 & 단가·계약 조건 검토, 폐기율 측정 담당 & 요구 제공자·측정자 & 계약 조건 명확 & 5,000만 원 이상 CFO 결재 지연 & 서면 \\ \hline
S-16 & 정하늘 · 감리법인 수석감리원 & 외부(계약) & 외부 감리 4.2억 수행 & 규제·감독 & RTM 추적성, 시험 완료 기준 문서 & 지적사항 미이행 & 서면 \\ \hline
S-17 & 타케다 마사시 · P2 설비 PM & 외부(계약, P2) & 설비 인터페이스 상대 & 공급자(간접) & 인터페이스 규격 조기 확정 & 규격 변경 & 서면. 리드타임 41주 통지 \\ \hline
S-18 & ㈜대성로지스 (3PL) & 외부(계약, 물류본부) & WMS 소유자 & 공급자(간접) & 계약 조건 유지 & 데이터 개방, 시스템 부하 & 2022 데이터 개방 거부 이력 \\ \hline
S-19 & 자사앱 외주 벤더 & 외부(계약, 영업본부) & 앱 소스 형상 보유 & 공급자 & 유지보수 계약 지속 & P1 API 연동 공수 & 월 3,200만 원 계약 \\ \hline
S-20 & 상용SW 벤더 3사 (API GW·MDM·APM) & 외부(계약) & 라이선스 공급 & 공급자 & 선약정 & 발주 지연(CFO 결재) & — \\ \hline
S-21 & 프로그램 PMO (4명) & 내부 & 프로그램 조정, 게이트 준비 & 의사결정 지원 & 프로젝트 간 정합 & P1·P2 일정 충돌 & — \\ \hline
S-22 & 재경팀장 & 내부 & CCB 위원, 원가 검증, 손실 확정 & 의사결정 지원 & 집행 한도 준수 & 예산 층 혼동 & — \\ \hline
S-23 & 규제기관 (개인정보보호위원회·금융위) & 외부(비계약) & 심사 주체 & 규제 & 법정 요건 충족 & — & — \\ \hline
S-24 & 직영 매장 직원 1,510명 / 가맹점주 331명 & 내부/외부 & 최종 사용자 & 영향받는 자 & 업무 부담 감소 & 전환기 이중 업무 & — \\ \hline
S-25 & 사모펀드 (2019년 350억 투자자) & 외부 & 이사회 관찰자, IPO 조항 & 영향받는 자·간접 의사결정 & 2027 IPO & 디지털 전환 미완 & — \\ \hline
\fi
\end{longtable}}

% ------------------------------------------------------------------ B
\secbar[70mm]{B. 등록부 — 현저성 3속성 · 태도 · 관여 전략 (항목 4~7)}
\guide{P(권력: \textbf{이 프로젝트를} 바꿀 수 있는가 — 직급이 아니다) / L(정당성: 요구가 조직·계약·법적으로 정당한가) / U(긴급성: 지금 대응을 요구하는가)를 상/중/하로 판정하고 Mitchell–Agle–Wood 7유형으로 조합. 태도는 무인지/저항/중립/지지/주도 “현재→목표”. 관여 전략(정보 제공/협의/공동 결정/협상/모니터링)+채널·빈도, 관계 소유자 실명(외부 협상 상대는 소관 임원). 속성이 바뀔 사건·게이트를 예측해 다음 검토 시점을 적는다.}
{\scriptsize
\begin{longtable}{|L{9mm}|L{12mm}|L{12mm}|L{12mm}|L{20mm}|L{24mm}|L{66mm}|L{24mm}|L{66mm}|}\hline
\hd{ID} & \hd{P} & \hd{L} & \hd{U} & \hd{현저성 유형} & \hd{태도 현재→목표} & \hd{관여 전략 · 채널 · 빈도} & \hd{관계 소유자} & \hd{속성 변화 예측 / 다음 검토} \\ \hline \endhead
\ifwbblank
\rd{12mm}S-01 & \g{상/중/하} & \g{상/중/하} & \g{상/중/하} & \g{[7유형 중]} & \g{[현재 → 목표]} & \g{[정보 제공/협의/공동 결정/협상/모니터링] + 채널·빈도} & \g{[실명]} & \g{[사건·게이트, 어느 유형으로]} \\ \hline
\rd{12mm}S-02 & & & & & & & & \\ \hline
\rd{12mm}S-03 & & & & & & & & \\ \hline
\rd{12mm}S-04 & & & & & & & & \\ \hline
\rd{12mm}S-05 & & & & & & & & \\ \hline
\rd{12mm}S-06 & & & & & & & & \\ \hline
\rd{12mm}S-07 & & & & & & & & \\ \hline
\rd{12mm}S-08 & & & & & & & & \\ \hline
\rd{12mm}S-09 & & & & & & & & \\ \hline
\rd{12mm}S-10 & & & & & & & & \\ \hline
\rd{12mm}S-11 & & & & & & & & \\ \hline
\rd{12mm}S-12 & & & & & & & & \\ \hline
\rd{12mm}S-13 & & & & & & & & \\ \hline
\rd{12mm}S-14 & & & & & & & & \\ \hline
\rd{12mm}S-15 & & & & & & & & \\ \hline
\rd{12mm}S-16 & & & & & & & & \\ \hline
\else
S-01 & 상 & 상 & 상 & \textbf{확정} & 지지(조건부)→주도 & 공동 결정. 주간 1:1 30분 + 월 CCB + 예외 보고 & P1 PM & 148억 축소 확정 시 긴급성 최고. G1 \\ \hline
S-02 & 상 & 상 & 중 & 지배 & 주도(일정)→주도(일정+범위 이연 수용) & 기대 관리. 게이트 보고 + 분기 이사회. 범위 이연 옵션을 미리 제시 & 정미란(스폰서 경유) & G3 전 긴급성 상승→확정. G2 \\ \hline
S-03 & 상 & 상 & 상 & \textbf{확정} & 저항(조건부)→지지 & 공동 결정. 격주 사용자 검토회, 프로토타입 우선 시연, 배포 창 설계 공동 확정 & P1 PM & 부속서 협상(2026-08) 전후 긴급성 최고. G1, 2026-08 \\ \hline
S-04 & 상 & 상 & 상 & \textbf{확정} & 중립(경계)→주도 & 공동 결정. 주간 아키텍처 리뷰. 인터페이스 13본 문서화를 “복원”으로 명명, 그의 지식을 공식 자산으로 & P1 PM & 데이터 전환 정본 판정(G2) 시 유일한 지식 보유자. G1 \\ \hline
S-05 & 상 & 중(P1에 대해) & 중 & 지배 & 중립→지지 & 협의. 월 1회 재고 원장 규격 검토, 3PL 협상 지원 자료 제공 & P1 PM (3PL 협상은 한동석 본인) & 3PL 협상(2026-06) 시 긴급성 상승. P2 지연 시 P1과 이해 충돌. G1 \\ \hline
S-06 & 중(CCB 표) & 상 & 하 & 재량→지배 가능 & 중립→지지 & 정보 제공+협의. 분기 1회 품질 데이터 요구 검토, 재고 정확도 분모 확장 협의 & P1 PM & 분모 확장 협의(M+3) 시 정당성 활용. G1 \\ \hline
S-07 & 중 & 상 & 상 & 의존→규제 기한 접근 시 확정 & 지지→지지 & 협의. 설계 스프린트마다 규제 체크포인트, 서면 검토 의견 기록 & P1 PM & 2026-09-11 전 8주 긴급성 최고. G1 \\ \hline
S-08 & 상(G4 판정권) & 상 & 하(현재) & 지배(G4에서 확정) & 중립→지지 & 협의. SAC 초안 공동 작성(G1까지), 월 1회 운영 요건 검토 & P1 PM & G3 이후 긴급성 최고. G1, G3 \\ \hline
S-09 & 상 & 상 & 상 & \textbf{확정} & 중립(방어)→지지 & 협상+공동 결정. 격주 변경협의체, 주간 PM 회의, 발주사 결정 기한 SLA 명문화 & P1 PM & 기준선 재확인(분기) 시. G1 \\ \hline
S-10 & 중→상(동의 거부권) & 상 & 중→상 & 의존→\textbf{확정}(2026-08) & 중립(분리 지지)→지지 & 협상(소유자 오정근). P1은 포털 프로토타입 분기 시연, 데이터 용도 명세서 제공 & 오정근 & 부속서 협상 2026-08-14 전후. 2026-06 \\ \hline
S-11 & 하 & 상 & 중 & 의존 & 저항(불신)→지지(슈퍼유저) & 협의. 설계 스프린트 사용자 검토 참여, 파일럿 12개점 후보 & P3 PM (P1은 검토회 초청) & 컷오버 준비(2027-01)에서 긴급성 상승. G2 \\ \hline
S-12 & 중(감사 착수 시 상) & 상 & 중 & 의존→감사 시 확정 & 중립→지지 & 정보 제공. 편익 표·컨틴전시 집행 기록 월 사본, 분기 서면 회신 & P1 PM & 감사 착수 시. G1 \\ \hline
S-13 & 상(현장 협조 거부) & 하(P1에 대해) & 상 & \textbf{위험} & 저항→중립 & 모니터링+정보 제공(소유자 한동석). P1은 “P1 편익에 물류 인건비 없음”을 문서화, 현장 시험 일정 사전 공유 & 한동석 (인사본부 협의) & P2 편익 논쟁이 P1로 번지면 정당성 획득→확정. G1 \\ \hline
S-14 & 중 & 상 & 하 & 재량→계약 시 지배 & 무인지→중립 & 정보 제공+협의. 별도 법인 계약·개인정보 처리 위탁 구조 G1까지 확정 & P1 PM (법무 최영주 지원) & 공장 연동 설계(2026-03). G1 \\ \hline
S-15 & 중 & 상 & 중 & 의존 & 중립→지지 & 협의. 조달 일정 공동 관리, 발주 리드타임 산정 & P1 PMO & 라이선스 발주(2026-02~03). G1 \\ \hline
S-16 & 상(검수 근거) & 상 & 하→시험기 상 & 지배 & 중립→중립 & 정보 제공. 감리 단계별 사전 협의, 지적 대응 기록 & P1 PMO & G2·G3 감리 시. G2 \\ \hline
S-17 & 중 & 중 & 상 & 요구→의존 & 무인지→중립 & 정보 제공. P2 PM 경유 규격 전달(T0 2026-04-24 전) & 프로그램 PMO / P2 PM & T0 전. G1 \\ \hline
S-18 & 상(연동 거부권) & 중 & 하 & 휴면→협상 시 위험 & 저항(이력)→중립 & 협상(소유자 한동석). P1은 연동 명세·부하 산정 제공 & 한동석 & 3PL 협상 2026-06. G1 \\ \hline
S-19 & 상(소스 보유) & 중 & 중 & 위험(소스 인질) & 중립→지지 & 협상(소유자 오정근·박세연). 소스 인수 조항 계약 변경, API 규격 조기 제공 & 박세연 & 앱 연동 설계(2026-04). G1 \\ \hline
S-20 & 중 & 상 & 중 & 의존 & 지지→지지 & 정보 제공. 조달 일정 공유 & P1 PMO & 발주 2026-02~03 \\ \hline
S-21 & 중 & 상 & 상 & 의존 & 지지→주도 & 공동 작업. 주간 프로그램 회의 & P1 PM & 상시 \\ \hline
S-22 & 중 & 상 & 중 & 의존 & 중립→지지 & 협의. 월 원가 보고 공동 검토 & P1 PMO & 월 \\ \hline
S-23 & 상 & 상 & 기한 접근 시 상 & 지배→확정 & 무인지→(해당 없음) & P4·CPO 경유. P1 직접 접촉 없음 & 최영주 (P4) & 2026-09, 2027-03 \\ \hline
S-24 & 하(개별)/상(집단) & 상 & 컷오버 시 상 & 의존→확정(컷오버) & 무인지→지지 & P3 경유 교육·파일럿. P1은 UAT 참여자 선정 지원 & P3 PM & 파일럿(2027-01). G2 \\ \hline
S-25 & 상(간접) & 상 & 하 & 지배(휴면에 가까움) & 무인지→(해당 없음) & 정미란·김선호 경유. P1 직접 접촉 없음 & 정미란 & IPO 준비(2027-H1) \\ \hline
\fi
\end{longtable}}

% ------------------------------------------------------------------ C
\secbar[95mm]{C. 현저성 분석 요약 \B{(판정 시점: \g{[게이트/날짜]})}{(G0 2026-01-05 기준)}}
\guide{7유형별로 해당 이해관계자를 모으고 유형별 대응 원칙을 적는다. 유형별 원칙은 미리 채워 두었다 — 자기 프로젝트에 맞게 고쳐 쓰라.}
{\footnotesize
\begin{tabularx}{\linewidth}{|L{30mm}|L{14mm}|Y|Y|}\hline
\hd{유형} & \hd{속성} & \hd{해당 이해관계자} & \hd{대응 원칙} \\ \hline
\ifwbblank
\rd{7mm}확정 (definitive) & P+L+U & & 최우선 관여, 공동 결정 \\ \hline
\rd{7mm}지배 (dominant) & P+L & & 기대 관리, 정기 협의 \\ \hline
\rd{7mm}위험 (dangerous) & P+U & & 정당성 경로 제공, 소관 임원이 관계 소유 \\ \hline
\rd{7mm}의존 (dependent) & L+U & & 권력 있는 후원자와 연결 \\ \hline
\rd{7mm}휴면 (dormant) & P & & 모니터링, 활성화 조건 기록 \\ \hline
\rd{7mm}재량 (discretionary) & L & & 정보 제공 \\ \hline
\rd{7mm}요구 (demanding) & U & & 경청, 에스컬레이션 경로 안내 \\ \hline
\else
확정 (definitive) & P+L+U & 정미란, 오정근, 박세연, 이재현 & 주간 접촉, 공동 결정. 이 넷의 합의 없이 기준선을 동결하지 않는다 \\ \hline
지배 (dominant) & P+L & 김선호, 한동석, 박준영, 정하늘, 규제기관, 사모펀드 & 기대 관리·정기 협의. 긴급성이 생기는 시점(G3, 3PL 협상, 감리, 규제 기한)을 미리 표시 \\ \hline
위험 (dangerous) & P+U & 김미르, 자사앱 벤더, (협상기의) 대성로지스 & 정당성 경로 제공(노조: 인사본부 협의 채널 / 앱 벤더: 소스 인수 계약 변경). 관계 소유자는 소관 임원 \\ \hline
의존 (dependent) & L+U & 서정필, 최영주, 강윤지, 윤태호, 강현도, 프로그램 PMO, 재경팀, 상용SW 벤더, 매장 직원·가맹점주 & 권력 있는 후원자와 연결(서정필↔오정근, 강윤지↔P3, 윤태호↔정미란) \\ \hline
휴면 (dormant) & P & 대성로지스(평시) & 활성화 조건(연동 요청) 기록, 협상 전 자료 준비 \\ \hline
재량 (discretionary) & L & 나영선, 조성태 & 정보 제공. 정당성이 P1 범위와 연결되는 시점(분모 확장, 공장 연동) 포착 \\ \hline
요구 (demanding) & U & 타케다(P2 설비 PM, 규격 요구) & 경청, P2 PM 경유 대응 \\ \hline
\fi
\end{tabularx}}

\begin{fieldbox}
\textbf{속성 변화 예측} \guide{(누가, 언제, 어느 유형으로 — 그 전에 무엇을 해 둘 것인가)}\par
\ifwbblank
\lines{3}{8mm}
\else
① 서정필: 2026-08-14 부속서 협상 시 388점 동의 거부권이 실질 권력이 되어 의존→확정. 협상 전 6월까지 포털 프로토타입 시연으로 태도를 지지로 옮겨 놓아야 한다. ② 김미르: P2 처리량·인력 논쟁이 P1 편익 표로 번지면 정당성을 획득해 위험→확정. P1 편익에 물류 인건비가 없음을 G1 전에 문서로 확정한다. ③ 박준영: G3 이후 SAC 판정권이 긴급해져 지배→확정. SAC 초안을 G1까지 공동 작성해 G4에서 놀라지 않게 한다. ④ 김선호: G3 go/no-go에서 확정. 범위 이연 옵션을 G2에서 미리 제시한다.
\fi
\end{fieldbox}



\end{document}
